% Generated draft: requires mathematical and source review.
\documentclass[11pt]{article}
\usepackage[margin=25mm]{geometry}
\usepackage{fontspec}
\setmainfont{Latin Modern Roman}
\usepackage{amsmath,amssymb,mathtools,unicode-math}
\setmathfont{Latin Modern Math}
\usepackage{enumitem,needspace,xurl,hyperref}
\setlength{\parindent}{0pt}
\setlength{\parskip}{5pt}
\setlength{\emergencystretch}{3em}
\newcommand{\N}{\mathbb{N}}
\newcommand{\Z}{\mathbb{Z}}
\newcommand{\Q}{\mathbb{Q}}
\newcommand{\R}{\mathbb{R}}
\newcommand{\C}{\mathbb{C}}
\newcommand{\F}{\mathbb{F}}
\newcommand{\A}{\mathbb{A}}
\newcommand{\PP}{\mathbb{P}}
\newcommand{\E}{\mathbb{E}}
\newcommand{\eps}{\varepsilon}
\newcommand{\dd}{\,\mathrm{d}}
\begin{document}
\section*{UnsolvedMath ID: 800974}
\textbf{Source problem code:} catalogue entry
\subsection{Definitions and mathematical statement}
Let \((M,g)\) be a compact Lorentzian manifold. A global conformal relation to a metric of constant curvature means that there exist a smooth Lorentzian metric \(h\) on \(M\), a smooth function \(u\colon M\to\R\), and a constant \(K\in\R\) such that
\[
g=e^{2u}h,
\]
and every nondegenerate tangent two-plane has sectional curvature \(K\) with respect to \(h\). A conformal diffeomorphism to another Lorentzian manifold of constant sectional curvature gives this formulation by pulling back its metric.

The question is whether these assumptions imply that \((M,g)\) is geodesically complete: does every maximal affinely parameterized \(g\)-geodesic \(\gamma\colon I\to M\) have \(I=\R\)? This asks for completeness of all geodesics, including timelike, null, and spacelike geodesics.

Here smooth manifolds without boundary and sectional curvature are the conventional interpretations of the source wording. The source does not explicitly specify a boundary convention, a dimension restriction, or regularity hypotheses; these interpretations require confirmation in mathematical review.
\subsection{Short English statement}
Must every compact Lorentzian manifold whose metric is globally conformal to a Lorentzian metric of constant sectional curvature be geodesically complete? Completeness concerns the original metric and affine parameters. The source supplies an open-problem label; the formulation and its conventions require human review.
\subsection{Sources}
\begin{itemize}
\item Local UnsolvedMath registry: permanent ID \(800974\), exact source problem code \texttt{catalogue entry}, in the AMR Open Problem Lists collection. The supplied record labels the problem \texttt{open}; this status is not independently verified.
\item Burns and Matveev, \emph{Open Problems and Questions about Geodesics}, source-list year supplied as 2019, Question 55. Supplied source URL: \url{https://arxiv.org/abs/1308.5417}. The record reports access on 2026-10-16 and extraction by \texttt{arxiv\_latex}; these are supplied provenance statements, not a claim of independent access.
\item The source attributes the surrounding results and list of open questions to the paper cited as \texttt{[Sanchez2015]}, from which this question is repeated. No full bibliographic entry or separate URL for that paper is supplied, so this attribution cannot be resolved further from the record.
\end{itemize}
\end{document}
